\IfFileExists{stacks-project.cls}{%
\documentclass{stacks-project}
}{%
\documentclass{amsart}
}

% For dealing with references we use the comment environment
\usepackage{verbatim}
\newenvironment{reference}{\comment}{\endcomment}
%\newenvironment{reference}{}{}
\newenvironment{slogan}{\comment}{\endcomment}
\newenvironment{history}{\comment}{\endcomment}

% For commutative diagrams we use Xy-pic
\usepackage[all]{xy}

% We use 2cell for 2-commutative diagrams.
\xyoption{2cell}
\UseAllTwocells

% We use multicol for the list of chapters between chapters
\usepackage{multicol}

% This is generally recommended for better output
\usepackage{lmodern}
\usepackage[T1]{fontenc}

% For cross-file-references
\usepackage{xr-hyper}

% Package for hypertext links:
\usepackage{hyperref}

% For any local file, say "hello.tex" you want to link to please
% use \externaldocument[hello-]{hello}
\externaldocument[introduction-]{introduction}
\externaldocument[conventions-]{conventions}
\externaldocument[sets-]{sets}
\externaldocument[categories-]{categories}
\externaldocument[topology-]{topology}
\externaldocument[sheaves-]{sheaves}
\externaldocument[sites-]{sites}
\externaldocument[stacks-]{stacks}
\externaldocument[fields-]{fields}
\externaldocument[algebra-]{algebra}
\externaldocument[brauer-]{brauer}
\externaldocument[homology-]{homology}
\externaldocument[derived-]{derived}
\externaldocument[simplicial-]{simplicial}
\externaldocument[more-algebra-]{more-algebra}
\externaldocument[smoothing-]{smoothing}
\externaldocument[modules-]{modules}
\externaldocument[sites-modules-]{sites-modules}
\externaldocument[injectives-]{injectives}
\externaldocument[cohomology-]{cohomology}
\externaldocument[sites-cohomology-]{sites-cohomology}
\externaldocument[dga-]{dga}
\externaldocument[dpa-]{dpa}
\externaldocument[sdga-]{sdga}
\externaldocument[hypercovering-]{hypercovering}
\externaldocument[schemes-]{schemes}
\externaldocument[constructions-]{constructions}
\externaldocument[properties-]{properties}
\externaldocument[morphisms-]{morphisms}
\externaldocument[coherent-]{coherent}
\externaldocument[divisors-]{divisors}
\externaldocument[limits-]{limits}
\externaldocument[varieties-]{varieties}
\externaldocument[topologies-]{topologies}
\externaldocument[descent-]{descent}
\externaldocument[perfect-]{perfect}
\externaldocument[more-morphisms-]{more-morphisms}
\externaldocument[flat-]{flat}
\externaldocument[groupoids-]{groupoids}
\externaldocument[more-groupoids-]{more-groupoids}
\externaldocument[etale-]{etale}
\externaldocument[chow-]{chow}
\externaldocument[intersection-]{intersection}
\externaldocument[pic-]{pic}
\externaldocument[weil-]{weil}
\externaldocument[adequate-]{adequate}
\externaldocument[dualizing-]{dualizing}
\externaldocument[duality-]{duality}
\externaldocument[discriminant-]{discriminant}
\externaldocument[derham-]{derham}
\externaldocument[local-cohomology-]{local-cohomology}
\externaldocument[algebraization-]{algebraization}
\externaldocument[curves-]{curves}
\externaldocument[resolve-]{resolve}
\externaldocument[models-]{models}
\externaldocument[functors-]{functors}
\externaldocument[equiv-]{equiv}
\externaldocument[pione-]{pione}
\externaldocument[etale-cohomology-]{etale-cohomology}
\externaldocument[proetale-]{proetale}
\externaldocument[relative-cycles-]{relative-cycles}
\externaldocument[more-etale-]{more-etale}
\externaldocument[trace-]{trace}
\externaldocument[crystalline-]{crystalline}
\externaldocument[spaces-]{spaces}
\externaldocument[spaces-properties-]{spaces-properties}
\externaldocument[spaces-morphisms-]{spaces-morphisms}
\externaldocument[decent-spaces-]{decent-spaces}
\externaldocument[spaces-cohomology-]{spaces-cohomology}
\externaldocument[spaces-limits-]{spaces-limits}
\externaldocument[spaces-divisors-]{spaces-divisors}
\externaldocument[spaces-over-fields-]{spaces-over-fields}
\externaldocument[spaces-topologies-]{spaces-topologies}
\externaldocument[spaces-descent-]{spaces-descent}
\externaldocument[spaces-perfect-]{spaces-perfect}
\externaldocument[spaces-more-morphisms-]{spaces-more-morphisms}
\externaldocument[spaces-flat-]{spaces-flat}
\externaldocument[spaces-groupoids-]{spaces-groupoids}
\externaldocument[spaces-more-groupoids-]{spaces-more-groupoids}
\externaldocument[bootstrap-]{bootstrap}
\externaldocument[spaces-pushouts-]{spaces-pushouts}
\externaldocument[spaces-chow-]{spaces-chow}
\externaldocument[groupoids-quotients-]{groupoids-quotients}
\externaldocument[spaces-more-cohomology-]{spaces-more-cohomology}
\externaldocument[spaces-simplicial-]{spaces-simplicial}
\externaldocument[spaces-duality-]{spaces-duality}
\externaldocument[formal-spaces-]{formal-spaces}
\externaldocument[restricted-]{restricted}
\externaldocument[spaces-resolve-]{spaces-resolve}
\externaldocument[formal-defos-]{formal-defos}
\externaldocument[defos-]{defos}
\externaldocument[cotangent-]{cotangent}
\externaldocument[examples-defos-]{examples-defos}
\externaldocument[algebraic-]{algebraic}
\externaldocument[examples-stacks-]{examples-stacks}
\externaldocument[stacks-sheaves-]{stacks-sheaves}
\externaldocument[criteria-]{criteria}
\externaldocument[artin-]{artin}
\externaldocument[quot-]{quot}
\externaldocument[stacks-properties-]{stacks-properties}
\externaldocument[stacks-morphisms-]{stacks-morphisms}
\externaldocument[stacks-limits-]{stacks-limits}
\externaldocument[stacks-cohomology-]{stacks-cohomology}
\externaldocument[stacks-perfect-]{stacks-perfect}
\externaldocument[stacks-introduction-]{stacks-introduction}
\externaldocument[stacks-more-morphisms-]{stacks-more-morphisms}
\externaldocument[stacks-geometry-]{stacks-geometry}
\externaldocument[moduli-]{moduli}
\externaldocument[moduli-curves-]{moduli-curves}
\externaldocument[examples-]{examples}
\externaldocument[exercises-]{exercises}
\externaldocument[guide-]{guide}
\externaldocument[desirables-]{desirables}
\externaldocument[coding-]{coding}
\externaldocument[obsolete-]{obsolete}
\externaldocument[fdl-]{fdl}
\externaldocument[index-]{index}

% Theorem environments.
%
\theoremstyle{plain}
\newtheorem{theorem}[subsection]{Theorem}
\newtheorem{proposition}[subsection]{Proposition}
\newtheorem{lemma}[subsection]{Lemma}

\theoremstyle{definition}
\newtheorem{definition}[subsection]{Definition}
\newtheorem{example}[subsection]{Example}
\newtheorem{exercise}[subsection]{Exercise}
\newtheorem{situation}[subsection]{Situation}

\theoremstyle{remark}
\newtheorem{remark}[subsection]{Remark}
\newtheorem{remarks}[subsection]{Remarks}

\numberwithin{equation}{subsection}

% Macros
%
\def\lim{\mathop{\mathrm{lim}}\nolimits}
\def\colim{\mathop{\mathrm{colim}}\nolimits}
\def\Spec{\mathop{\mathrm{Spec}}}
\def\Hom{\mathop{\mathrm{Hom}}\nolimits}
\def\Ext{\mathop{\mathrm{Ext}}\nolimits}
\def\SheafHom{\mathop{\mathcal{H}\!\mathit{om}}\nolimits}
\def\SheafExt{\mathop{\mathcal{E}\!\mathit{xt}}\nolimits}
\def\Sch{\mathit{Sch}}
\def\Mor{\mathop{\mathrm{Mor}}\nolimits}
\def\Ob{\mathop{\mathrm{Ob}}\nolimits}
\def\Sh{\mathop{\mathit{Sh}}\nolimits}
\def\NL{\mathop{N\!L}\nolimits}
\def\CH{\mathop{\mathrm{CH}}\nolimits}
\def\proetale{{pro\text{-}\acute{e}tale}}
\def\etale{{\acute{e}tale}}
\def\QCoh{\mathit{QCoh}}
\def\Ker{\mathop{\mathrm{Ker}}}
\def\Im{\mathop{\mathrm{Im}}}
\def\Coker{\mathop{\mathrm{Coker}}}
\def\Coim{\mathop{\mathrm{Coim}}}

% Boxtimes
%
\DeclareMathSymbol{\boxtimes}{\mathbin}{AMSa}{"02}

%
% Macros for moduli stacks/spaces
%
\def\QCohstack{\mathcal{QC}\!\mathit{oh}}
\def\Cohstack{\mathcal{C}\!\mathit{oh}}
\def\Spacesstack{\mathcal{S}\!\mathit{paces}}
\def\Quotfunctor{\mathrm{Quot}}
\def\Hilbfunctor{\mathrm{Hilb}}
\def\Curvesstack{\mathcal{C}\!\mathit{urves}}
\def\Polarizedstack{\mathcal{P}\!\mathit{olarized}}
\def\Complexesstack{\mathcal{C}\!\mathit{omplexes}}
% \Pic is the operator that assigns to X its picard group, usage \Pic(X)
% \Picardstack_{X/B} denotes the Picard stack of X over B
% \Picardfunctor_{X/B} denotes the Picard functor of X over B
\def\Pic{\mathop{\mathrm{Pic}}\nolimits}
\def\Picardstack{\mathcal{P}\!\mathit{ic}}
\def\Picardfunctor{\mathrm{Pic}}
\def\Deformationcategory{\mathcal{D}\!\mathit{ef}}
\usepackage{cleveref}
\crefname{originaldefinition}{original definition}{original definitions}
\Crefname{originaldefinition}{Original definition}{Original definitions}
\makeatletter
\newlabel{definition-separated}{{007A}{source}{}{sourcepointer-1}{}}
\newlabel{definition-separated@cref}{{[originaldefinition][0][]007A}{[1][0][]source}{}{}{}}
\makeatother
\crefname{originallemma}{original lemma}{original lemmas}
\Crefname{originallemma}{Original lemma}{Original lemmas}
\makeatletter
\newlabel{lemma-separated-presheaf-into-sheaf}{{0082}{source}{}{sourcepointer-2}{}}
\newlabel{lemma-separated-presheaf-into-sheaf@cref}{{[originallemma][0][]0082}{[1][0][]source}{}{}{}}
\makeatother
\crefname{originalsection}{original section}{original sections}
\Crefname{originalsection}{Original section}{Original sections}
\makeatletter
\newlabel{categories-section-posets-limits}{{002Z}{source}{}{sourcepointer-3}{}}
\newlabel{categories-section-posets-limits@cref}{{[originalsection][0][]002Z}{[1][0][]source}{}{}{}}
\makeatother
\title{Sections of directed colimits of sheaves versus colimits of sections}
\author{The Stacks Project Authors}
\date{stacks-project:sheaves.tex:local-snapshot-20261012; GNU FDL 1.2 or later}
\begin{document}
\maketitle
\setcounter{subsection}{1}
\noindent\textit{Editorial note.} All delivered mathematical text is the author's own, copied from the pinned Stacks Project chapter file sheaves.tex; no statement or proof step was rewritten, reordered or removed. The counted span is the author environment and its own proof environment. References inside the span are handed to the assembler as navigation stubs that carry the canonical tag of the cited tag only; whether each of them is load-bearing is recorded on the stub itself.\par\medskip
\noindent\textit{Allowed standard prerequisites (editorial).}\begin{itemize}
\item Definition definition-separated (canonical tag 007A): a presheaf is separated if restrictions agreeing on a covering come from a common section, i.e. the restriction maps to the sheafification are injective. invoked at proof lines 3688-3690 for the presheaf V \textbackslash{}mapsto colim F\_i(V) under injective transitions. Makes F' separated so the map to its sheafification is injective.
\item Lemma lemma-separated-presheaf-into-sheaf (canonical tag 0082): the canonical map from a separated presheaf into its sheafification is injective. cited at proof lines 3692-3694. Completes the proof of part (1).
\end{itemize}\medskip
\ifdefined\hypertarget\hypertarget{sourcepointer-1}{}\fi
\noindent\textit{Source reference.} Original definition 007A: sources/stacks-chow/tags.txt, source lines 274--274. Invoked at proof lines 3688-3690 for the separatedness of the colimit presheaf; load-bearing inside the counted span.\par\medskip
\ifdefined\hypertarget\hypertarget{sourcepointer-2}{}\fi
\noindent\textit{Source reference.} Original lemma 0082: sources/stacks-chow/tags.txt, source lines 302--302. Cited at proof lines 3692-3694 to conclude part (1); load-bearing inside the counted span.\par\medskip
\ifdefined\hypertarget\hypertarget{sourcepointer-3}{}\fi
\noindent\textit{Source reference.} Original section 002Z: sources/stacks-chow/tags.txt, source lines 116--116. Cited at statement lines 3666-3667 fixing the notion of a system of sheaves over the directed set I; navigation for the setting.\par\medskip
\begin{lemma}
\label{lemma-directed-colimits-sections}
Let $X$ be a topological space. Let $I$ be a directed set.
Let $(\mathcal{F}_i, \varphi_{ii'})$ be a system of sheaves of sets
over $I$, see
Categories, Section \ref{categories-section-posets-limits}.
Let $U \subset X$ be an open subset.
Consider the canonical map
$$
\Psi :
\colim_i \mathcal{F}_i(U)
\longrightarrow
\left(\colim_i \mathcal{F}_i\right)(U)
$$
\begin{enumerate}
\item If all the transition maps are injective then
$\Psi$ is injective for any open $U$.
\item If $U$ is quasi-compact, then $\Psi$ is injective.
\item If $U$ is quasi-compact and all the transition maps are injective
then $\Psi$ is an isomorphism.
\item If $U$ has a cofinal system of open coverings
$\mathcal{U} : U = \bigcup_{j\in J} U_j$ with
$J$ finite and $U_j \cap U_{j'}$ quasi-compact
for all $j, j' \in J$, then $\Psi$ is bijective.
\end{enumerate}
\end{lemma}

\begin{proof}
Assume all the transition maps are injective. In this case the presheaf
$\mathcal{F}' : V \mapsto \colim_i \mathcal{F}_i(V)$ is
separated (see Definition \ref{definition-separated}).
By the discussion above we have
$(\mathcal{F}')^\# = \colim_i \mathcal{F}_i$.
By Lemma \ref{lemma-separated-presheaf-into-sheaf} we see that
$\mathcal{F}' \to (\mathcal{F}')^\#$ is injective. This proves (1).

\medskip\noindent
Assume $U$ is quasi-compact.
Suppose that $s \in \mathcal{F}_i(U)$ and
$s' \in \mathcal{F}_{i'}(U)$ give rise to elements on
the left hand side which have the same image under $\Psi$.
Since $U$ is quasi-compact this means there exists
a finite open covering $U = \bigcup_{j = 1, \ldots, m} U_j$
and for each $j$ an index $i_j \in I$, $i_j \geq i$, $i_j \geq i'$
such that $\varphi_{ii_j}(s)$ and $\varphi_{i'i_j}(s')$ have the
same restriction to $U_j$.
Let $i''\in I$ be $\geq$ than all of the $i_j$.
We conclude that $\varphi_{ii''}(s)$ and $\varphi_{i'i''}(s')$
agree on the opens $U_j$ for all $j$ and hence that
$\varphi_{ii''}(s) = \varphi_{i'i''}(s')$. This proves (2).

\medskip\noindent
Assume $U$ is quasi-compact and all transition maps injective.
Let $s$ be an element of the target of $\Psi$.
Since $U$ is quasi-compact
there exists a finite open covering $U = \bigcup_{j = 1, \ldots, m} U_j$,
for each $j$ an index $i_j \in I$ and $s_j \in \mathcal{F}_{i_j}(U_j)$
such that $s|_{U_j}$ comes from $s_j$ for all $j$.
Pick $i \in I$ which is $\geq$ than all of the $i_j$.
By (1) the sections $\varphi_{i_ji}(s_j)$ agree over the
overlaps $U_j \cap U_{j'}$. Hence they glue to a section
$s' \in \mathcal{F}_i(U)$ which maps to $s$ under $\Psi$.
This proves (3).

\medskip\noindent
Assume the hypothesis of (4). In particular we see that
$U$ is quasi-compact and hence by (2) we have injectivity of $\Psi$.
Let $s$ be an element of the target of $\Psi$.
By assumption there  exists a finite open covering
$U = \bigcup_{j = 1, \ldots, m} U_j$, with $U_j \cap U_{j'}$
quasi-compact for all $j, j' \in J$ and
for each $j$ an index $i_j \in I$ and $s_j \in \mathcal{F}_{i_j}(U_j)$
such that $s|_{U_j}$ is the image of $s_j$ for all $j$.
Since $U_j \cap U_{j'}$ is quasi-compact we can apply (2)
and we see that there exists an $i_{jj'} \in I$,
$i_{jj'} \geq i_j$, $i_{jj'} \geq i_{j'}$ such that
$\varphi_{i_ji_{jj'}}(s_j)$ and $\varphi_{i_{j'}i_{jj'}}(s_{j'})$
agree over $U_j \cap U_{j'}$. Choose an index $i \in I$
which is bigger or equal than all the $i_{jj'}$. Then we see that
the sections $\varphi_{i_ji}(s_j)$ of $\mathcal{F}_i$ glue
to a section of $\mathcal{F}_i$ over $U$. This section is mapped
to the element $s$ as desired.
\end{proof}

\end{document}
